\section{Aula 1}

O Augusto fez as definições básicas de percolação
e apresentou resultados iniciais bem bonitos.

Dado um grafo $G = (V,E)$ definimos (assim como em $G(n,p)$)
o grafo aleatório $G_p$ que tem cada aresta adicionada
independentemente com probabilidade $p$. Aqui $V$ e $E$ são usualmente
infinitos. O exemplo canônico é $G = (\Z^d, \{\{x,y\} : |x - y| = 1\})$,
e percolamos sobre o látice inteiro de dimensão $d$.

O primeiro teorema foi bem simples.
\begin{theorem}
	Seja $G = (V,E)$ com $|V| = \infty$ e suponha que $\Delta(G) < \infty$. Então
	$$\Pr(G_p \text{ é desconexo}) = 1.$$
\end{theorem}
\begin{proof}
	Basta verificar que com probabilidade $q>0$ cada vtx é isolado.
	Como vertices não adjacentes são isolados independentemente, segue
	de Borel-Cantelli que
	isso acontece infinitas vezes com probabilidade 1.
\end{proof}

O Augusto motivou percolação para responder a seguinte pergunta
\begin{question}
	Qual é a probabilidade de que $G_p$ (nesse caso $\Z^d$)
	possua uma componente conexa infinita?
\end{question}

Com um lema esperto, podemos reduzir essa pergunta a probabilidade da origem
estar conectada ao infinito. Para ser mais exato, definimos o evento
\[
	\{\vec{0} \leftrightarrow \infty\} = \{\exists \sigma = (v_0, v_1,
	v_2, \dots) \subset G_p\}
\]
onde $\sigma$ aqui é um caminho infinito que não se intersecta e $v_0
= 0$. Temos então
\begin{lemma}
	\[
		\Pr(G_p \text{ tem uma comp. con. infinita}) =
		\begin{cases}
			1 & \text{se } \Pr_p(\vec{0} \leftrightarrow \infty) > 0\\
			0 & \text{se } \Pr_p(\vec{0} \leftrightarrow \infty) = 0
		\end{cases}
	\]
\end{lemma}
\begin{proof}
	Um lado é fácil, se a probabilidade de se conectar ao infinito é
	$0$, então, pelo
	union bound, nenhum vertice se conecta ao infinito e $G_p$ não tem
	componente infinita.

	Para o outro lado, basta notar que ter uma componente infinita é um
	evento caudal,
	portanto sua probabilidade é 0 ou 1. Se a probabilidade do
	$\vec{0}$ se conectar
	ao infinito é maior que $0$, então temos uma componente com probabilidade 1.
\end{proof}

Provamos outro lema esperto também.
\begin{lemma}
	\[
		\{\vec{0} \leftrightarrow \infty\} = \bigcap_{n\geq1} \{\exists
		\sigma_k = (\vec{0}, v_1, \dots, v_k) \subset G_p, v_k \in \partial B_n\}
	\]
	onde $\sigma_k$ é um caminho não intersectante com $k$
	arestas conectando o $0$ com a borda da bola quadrada de
	raio $n$.
\end{lemma}

\begin{proof}
	A primeira inclusão é óbvia, segue de casas dos pombos. A outra nem tanto,
	note que cada evento nos diz que existe um caminho, não
	necessáriamente esse caminho
	é o mesmo. Se ele fica saltando não adianta de nada. Mas é fácil
	resolver. Como existem
	infinitos caminhos que saem do $\vec{0}$, basta anotar uma aresta
	vizinha do $\vec{0}$ pela qual passam
	infinitos caminhos. Andamos pela a aresta e continuamos, novamente
	existe uma aresta (sem ser a anterior) pela
	qual passam infinitos caminhos. Percorrendo assim, terminamos a prova.
\end{proof}

Por fim o Augusto enunciou um teoremão bem bonito. Para reduzir a notação,
seja $\Theta_d(p) = \Pr_p(\vec{0} \leftrightarrow \infty)$ em $\Z^d$.
Temos então
\begin{theorem}
	\label{trm:bounds_for_pc}
	Se $p < \frac{1}{2d - 1}$, então
	\[\Theta(p) = 0.\]
	Por outro lado, se $p > 2/3$, então
	\[\Theta(p) > 0.\]
\end{theorem}
Ele só provou a parte fácil, que é a primeira afirmação.
\begin{proof}
	Basta usar o lema anterior e cotar a probabilidade de conectar o $0$
	com a borda de $B_n$.
	Existem menos que $(2d)(2d - 1)^{n-1}$ caminhos não intersectantes
	saindo do $0$ de comprimento $n$,
	cada um deles acontece com probabilidade $p^n$. Portanto,
	\[
		\Pr_p(\exists
			\sigma_k = (\vec{0}, v_1, \dots, v_k) \subset G_p, v_k \in \partial
		B_n) \leq (2d)(2d-1)^{n-1}p^n \to 0
	\]
	se $p < \frac{1}{2d - 1}$.
\end{proof}

Próxima aula veremos a segunda parte do teorema (parece bem interessante).
